\BeleskeID{09_2-s025}
% element: 09_2-s025:e04
Daljim povećanjem ulaznog napona izlazni napon opada, dok napon između sorsa i drejna P tranzistora raste. Oblast svakog tranzistora određuje se poređenjem napona drejn–sors sa granicom zasićenja njegovog modela. U oznakama priključaka ti uslovi su
\[V_{DS,N}\ge\frac{(V_{GS,N}-V_{Tn})}{1+\frac{(V_{GS,N}-V_{Tn})}{L_nE_{Cn}}}\]
\[V_{DS,P}\ge\frac{(V_{GS,P}-V_{Tp})}{1+\frac{(V_{GS,P}-V_{Tp})}{L_pE_{Cp}}}\]
N tek počinje da vodi malu struju pri $V_I=V_{Tn}+\varepsilon$. Visok napon drejna daje zasićenje, a mali napon sors–drejn P tranzistora omsku oblast.
